\par\Needspace{6\baselineskip}\medskip
\hypertarget{native-084-article065}{}\pdfbookmark[2]{65.}{article-065}
\begin{center}65.\end{center}
When only a very few geocentric positions are to be determined from given elements, it is scarcely worth using all the devices described above that allow us to pass directly from the eccentric anomaly to geocentric longitude and latitude, or even to right ascension and declination. The resulting savings would be swallowed up by the number of auxiliary quantities that must first be calculated. Yet combining the reduction to the ecliptic with the calculation of geocentric longitude and latitude will always yield an appreciable saving. For if the ecliptic itself is taken as the plane for coordinate $z$, and the poles of the planes for coordinates $x$, $y$ are placed at longitudes $\AscendingNode$, $90^0+\AscendingNode$, the coordinates are very easily determined without the need for any auxiliary quantities. In fact, we have
\SourceLeaf{085}
\[
\begin{array}{@{}l|l|l@{}}
x=r\cos u & X=R'\cos(L-\AscendingNode) & x-X=\Delta'\cos(l-\AscendingNode)\\[.4em]
y=r\cos i\sin u & Y=R'\sin(L-\AscendingNode) & y-Y=\Delta'\sin(l-\AscendingNode)\\[.4em]
z=r\sin i\sin u & Z=R'\operatorname{tang}B & z-Z=\Delta'\operatorname{tang}b.
\end{array}
\]
\Needspace{12\baselineskip}
\noindent Whenever $B=0$, we have $R'=R$, $Z=0$. Using these formulas, we work out our example with the following numbers, starting from $L-\AscendingNode=213^0\,12^{\prime}\,0^{\prime\prime}{,}32$
\begin{center}
\begin{tabular}{@{}ll@{}}
\makebox[13em][l]{$\log r$\hfill .\,.\,.\,.\,.\,.\,.\,.\,.\,.\,.} & $0{,}325\,9877$\\[.3em]
\makebox[13em][l]{$\log\cos u$\hfill .\,.\,.\,.\,.\,.\,.\,.} & $9{,}982\,4141_n$\\[.3em]
\makebox[13em][l]{$\log\sin u$\hfill .\,.\,.\,.\,.\,.\,.\,.} & $9{,}445\,4714_n$\\[.3em]
\hline
\makebox[13em][l]{$\log x$\hfill .\,.\,.\,.\,.\,.\,.\,.\,.\,.\,.} & $0{,}308\,4018_n$\\[.3em]
\hline
\makebox[13em][l]{$\log r\sin u$\hfill .\,.\,.\,.\,.\,.\,.} & $9{,}771\,4591_n$\\[.3em]
\makebox[13em][l]{$\log\cos i$\hfill .\,.\,.\,.\,.\,.\,.\,.} & $9{,}988\,5266$\\[.3em]
\makebox[13em][l]{$\log\sin i$\hfill .\,.\,.\,.\,.\,.\,.\,.} & $9{,}355\,7570$\\[.3em]
\hline
\makebox[13em][l]{$\log y$\hfill .\,.\,.\,.\,.\,.\,.\,.\,.\,.\,.} & $9{,}759\,9857_n$\\[.3em]
\makebox[13em][l]{$\log z$\hfill .\,.\,.\,.\,.\,.\,.\,.\,.\,.\,.} & $9{,}127\,2161_n$
\end{tabular}
\end{center}
\begin{center}
\begin{tabular}{@{}ll@{}}
\makebox[13em][l]{$\log R'$\hfill .\,.\,.\,.\,.\,.\,.\,.\,.\,.} & $9{,}998\,0979$\\[.3em]
\makebox[13em][l]{$\log\cos(L-\AscendingNode)$\hfill .\,.\,.} & $9{,}922\,6027_n$\\[.3em]
\makebox[13em][l]{$\log\sin(L-\AscendingNode)$\hfill .\,.\,.} & $9{,}738\,4353_n$\\[.3em]
\hline
\makebox[13em][l]{$\log X$\hfill .\,.\,.\,.\,.\,.\,.\,.\,.\,.} & $9{,}920\,7006_n$\\[.3em]
\makebox[13em][l]{$\log Y$\hfill .\,.\,.\,.\,.\,.\,.\,.\,.\,.} & $9{,}736\,5332_n$\\[.3em]
$Z=$ & $0$
\end{tabular}
\end{center}
\noindent This gives
\begin{center}
\begin{tabular}{@{}ll@{}}
\makebox[13em][l]{$\log(x-X)$\hfill .\,.\,.\,.\,.\,.} & $0{,}079\,5906_n$\\[.3em]
\makebox[13em][l]{$\log(y-Y)$\hfill .\,.\,.\,.\,.\,.} & $8{,}480\,7165_n$\\[.3em]
whence $(l-\AscendingNode)=$ & $181^0\,26^{\prime}\,33^{\prime\prime}{,}49$\\[.3em]
\makebox[13em][l]{$\log\Delta'$\hfill .\,.\,.\,.\,.\,.\,.\,.\,.\,.} & $0{,}079\,7283$\\[.3em]
\makebox[13em][l]{$\log\operatorname{tang}b$\hfill .\,.\,.\,.\,.\,.\,.} & $9{,}047\,4878_n$
\end{tabular}
\end{center}
\begin{center}
\begin{tabular}{@{}rl@{}}
$l=$ & $352^0\,34^{\prime}\,22^{\prime\prime}{,}22$\\[.4em]
$b=$ & $-6\quad21\quad55{,}06$
\end{tabular}
\end{center}
